%!TEX root = ../manuscript.tex

\begin{theorem}[Stuck-State Characterization, permissive version]
\label{thm:stuck-permissive}
Let $w$ be a non-empty Gauss word and let $N(w)$ be its nesting
forest. Then $w$ is stuck under \emph{permissive} R1-down and R2-down
moves if and only if both of the following hold:
\begin{enumerate}[(i)]
\item no adjacent opposite-sign pair exists in $w$ (no R1-down
  applicable); and
\item $N(w)$ has no edges, i.e.\ no chord directly nests any other
  chord.
\end{enumerate}
\end{theorem}

\begin{proof}
($\Rightarrow$) We prove the contrapositive. If an R1-down move is
applicable, then $w$ is not stuck. If $N(w)$ has an edge, then some
chord $a$ directly nests some chord $b$, and permissive R2-down on
$(a, b)$ is available (the permissive move requires only direct
nesting), so $w$ is not stuck.

($\Leftarrow$) Suppose both (i) and (ii) hold. By (i), no R1-down is
applicable. By (ii), no chord directly nests any other, so no
permissive R2-down is applicable. Hence $w$ is stuck under permissive
moves.
\end{proof}

\begin{theorem}[Stuck-State Characterization, bigon version]
\label{thm:stuck-bigon}
Let $w$ be a non-empty Gauss word and let $N(w)$ be its nesting
forest. Then $w$ is stuck under \emph{bigon} R1-down and R2-down
moves if and only if both of the following hold:
\begin{enumerate}[(i)]
\item no adjacent opposite-sign pair exists in $w$ (no R1-down
  applicable); and
\item no edge $(a, b)$ of $N(w)$ satisfies the bigon condition
  (both outer arcs of $(a, b)$ are empty of other chord endpoints).
\end{enumerate}
\end{theorem}

\begin{proof}
($\Rightarrow$) Contrapositive. If an R1-down is applicable, not
stuck. If some directly-nested pair $(a, b) \in N(w)$ satisfies the
bigon condition, then bigon R2-down on $(a, b)$ is applicable, so not
stuck.

($\Leftarrow$) By (i), no R1-down. By (ii), every directly-nested
pair fails the bigon condition (at least one outer arc contains a
chord endpoint), so no bigon R2-down is applicable. Moreover, a
bigon R2-down pair must be directly nested (the bigon condition for
a nested but not directly-nested pair $(a, c)$ requires the outer
arcs to be empty, but the intermediate chord $b$ with $a \succ b \succ
c$ has at least one endpoint in those outer arcs, contradicting the
condition). Hence no bigon R2-down is applicable at all, and $w$ is
stuck.
\end{proof}

\begin{corollary}[Relation to interlacement structure]
\label{cor:interlace}
A Gauss word $w$ is stuck under permissive R1/R2 moves if and only if
every pair of distinct chords in $w$ is either interleaved or
disjoint.
\end{corollary}

\begin{proof}
Two chords $a, b$ are related in exactly one of three ways: $a$
nests $b$, $b$ nests $a$, or $a, b$ are interleaved or disjoint (not
nested). ``No pair is nested'' is therefore equivalent to ``every
pair is interleaved or disjoint.''

``$N(w)$ has no edges'' (condition (ii) of Theorem~\ref{thm:stuck-permissive})
is equivalent to ``no chord directly nests another.'' If some chord
$a$ nests some chord $b$ but not directly, by Lemma~\ref{lem:E} there
is a chain of direct nestings $a = c_0 \succ c_1 \succ \cdots \succ
c_k = b$ with $k \geq 2$ in $N(w)$, giving at least $k \geq 1$ edge
of $N(w)$. Hence ``no nesting'' is equivalent to ``no direct
nesting,'' completing the proof.
\end{proof}

\begin{remark}
Corollary~\ref{cor:interlace} is strictly weaker than the stronger
statement ``every interlacement component of $w$ is a fully
interleaved clique,'' which was used in an earlier draft. The weaker
form is correct because disjoint chords can coexist within a single
interlacement component when interleaving is transitive but not
complete (e.g., a chord $a$ interleaved with $b$, $b$ interleaved
with $c$, yet $a$ and $c$ disjoint).
\end{remark}
